\documentclass[conference]{IEEEtran}

\usepackage[T1]{fontenc}
\usepackage[utf8]{inputenc}
\usepackage{cite}
\usepackage{amsmath}
\usepackage{graphicx}
\usepackage{booktabs}
\usepackage{hyperref}

\title{Ethernet-APL-Based Cloud--Edge Collaborative Architecture\\
       for Industrial Automation Systems}

\author{
  \IEEEauthorblockN{Author One}
  \IEEEauthorblockA{
    Hochschule Hamm-Lippstadt\\
    Hamm, Germany\\
    author1@hshl.de
  }
  \and
  \IEEEauthorblockN{Author Two}
  \IEEEauthorblockA{
    Hochschule Hamm-Lippstadt\\
    Hamm, Germany\\
    author2@hshl.de
  }
  \and
  \IEEEauthorblockN{Author Three}
  \IEEEauthorblockA{
    Hochschule Hamm-Lippstadt\\
    Hamm, Germany\\
    author3@hshl.de
  }
  \and
  \IEEEauthorblockN{Author Four}
  \IEEEauthorblockA{
    Hochschule Hamm-Lippstadt\\
    Hamm, Germany\\
    author4@hshl.de
  }
}

\begin{document}
\maketitle

% ----------------------------------------------------------
% Abstract + Introduction (~1 page)
% ----------------------------------------------------------
\begin{abstract}

\end{abstract}

\begin{IEEEkeywords}
Ethernet-APL, cloud--edge collaborative, industrial automation,
quality of service, time-sensitive networking, OPC UA
\end{IEEEkeywords}

\section{Introduction}\label{sec:intro}

Example citation \cite{placeholder2025}.

\subsection{Motivation}

\subsection{Problem Statement}

\subsection{Contributions}

% ----------------------------------------------------------
% Literature Review (~1 page)
% ----------------------------------------------------------
\section{Literature Review}\label{sec:literature}

\subsection{Evolution of Fieldbus and Industrial Ethernet}

\subsection{Ethernet-APL Standard}

\subsection{Cloud--Edge Collaboration}

\subsection{TSN and OPC UA}

% ----------------------------------------------------------
% Proposed Method (~1.5 pages)
% ----------------------------------------------------------
\section{Proposed Architecture}\label{sec:method}

\subsection{System Architecture Overview}

\begin{figure}[!t]
  \centering
  % \includegraphics[width=\columnwidth]{figures/architecture.pdf}
  \rule{\columnwidth}{0.4pt}
  \caption{Ethernet-APL-based cloud--edge collaborative architecture.}
  \label{fig:architecture}
\end{figure}

\subsection{Field Layer}

\subsection{Edge Layer}

\subsection{Cloud Layer}

\subsection{Data Flow Model}

\subsection{QoS Scheduling}

% ----------------------------------------------------------
% Evaluation (~1.5 pages)
% ----------------------------------------------------------
\section{Evaluation}\label{sec:evaluation}

\subsection{Experimental Setup}

\subsection{Application Example}

\subsection{Communication Requirements}

\subsection{Performance Metrics: Latency, Throughput, Jitter}

\subsection{Pros and Cons of Ethernet-APL}

% ----------------------------------------------------------
% Conclusion (~1 page incl. references)
% ----------------------------------------------------------
\section{Conclusion}\label{sec:conclusion}

\bibliographystyle{IEEEtran}
\bibliography{references}

\end{document}
